\chapter{Desarrollo de la comparativa: modelos y protocolo}
\label{cap:comparativa}

% GUION. Fusion de los antiguos capitulos 8 (Los tres modelos) y 9 (Protocolo
% experimental). Corresponde al capitulo 5 de la estructura del tipo 3:
% «Desarrollo de la comparativa — con todo detalle, con todos los resultados y
%  mediciones obtenidos».
%
% Fuentes: TFM-Quantum/src/features.py, src/particion.py, doc/plan_eda.md,
%          configs/gem_config.yaml, doc/memoria/materiales.md seccion 5.
% ATENCION: el esqueleto anterior describia un modulo REM por GNN que esta FUERA DE
% ALCANCE, y una validacion en hardware Eagle/Osprey que tampoco lo esta —Eagle esta
% retirada y el plan abierto da 10 min/mes—. Si se mencionan, que sea como trabajo futuro.

\section{El problema de la representación}
\label{sec:representacion_modelos}
% ESTA seccion es el nucleo conceptual de la comparativa: es lo que hace que comparar
% estos tres modelos y no otros tres tenga sentido.
Cada muestra es una \textbf{matriz de altura variable} (de una decena a varios miles de
filas, una por puerta). Una \textbf{red de grafos} acepta esa entrada tal cual. Ridge y
Random Forest \textbf{no}: exigen que todas las filas tengan las mismas columnas.

\subsection{La agregación, y qué se pierde con ella}
Resumir el circuito en una fila \textbf{tira el orden y la estructura}: dos circuitos con las
mismas puertas en distinto orden dan \textbf{la misma fila}, y son físicamente distintos.

\textbf{Ahí está el núcleo de este trabajo:} la comparativa mide exactamente cuánto vale esa
información que el grafo conserva y la tabla pierde.

\subsection{Las variables agregadas}
\label{sec:features}
% Fuente: src/features.py (modulo de produccion, catalogo cerrado)
Ocho familias. Describir cada una en una frase y detallar solo las que resultaron
predictivas:
\begin{enumerate}
    \item tamaño y forma del circuito;
    \item composición en tipos de puerta;
    \item ángulos de rotación;
    \item calidad de las puertas ejecutadas;
    \item telemetría del hardware que le tocó;
    \item cocientes físicos: duración frente a decoherencia;
    \item la heurística del estado del arte (el score de mapomatic como una variable más);
    \item tiempo.
\end{enumerate}

\subsection{Dos decisiones de diseño que importan}
% Fuente: src/features.py, docstrings de _telemetria_por_uso y _profundidad_y_camino
\textbf{Telemetría ponderada por uso:} un qubit malo que interviene en $200$ puertas debe pesar
más que uno bueno usado dos veces.

\textbf{Duración crítica, no suma de duraciones:} las puertas sobre qubits distintos se
ejecutan simultáneamente. Lo que vive el estado es el \textbf{camino crítico del DAG}.

\subsection{Qué se excluye a propósito}
El \textbf{tipo de circuito} no es una variable: en entrenamiento es constante, y en test le
daría al modelo la respuesta al eje que precisamente se quiere medir.

% =====================================================================================
% LAS TRES SOLUCIONES COMPARADAS
% =====================================================================================

\section{Solución 1: Ridge}
\label{sec:ridge}
\subsection{Por qué Ridge y no mínimos cuadrados}
% Fuente: doc/plan_eda.md seccion 4. Cifras del v1: REHACER sobre el v2.
Medido: la colinealidad entre variables es \textbf{extrema} (un bloque entero es «el tamaño»
escrito de ocho maneras). Con columnas casi idénticas, mínimos cuadrados no puede decidir a
cuál atribuir el mérito: produce coeficientes enormes que \textbf{solo se cancelan dentro del
rango visto}. Y nuestro test es OOD en tres ejes.

\section{Solución 2: Random Forest}
\label{sec:rf}
El mejor baseline tabular según Liao et al. (2024). \textbf{Y sirve para algo más:} es inmune
a la colinealidad (cada árbol elige una columna y sigue), así que la diferencia con Ridge no
es un problema a resolver, sino \textbf{parte de lo que la comparativa mide}.

\textbf{Su limitación, que hay que declarar:} \textbf{no puede extrapolar}. Satura en el valor
de la hoja más extrema que aprendió, y este conjunto de datos activa ese caso a propósito.

\section{Solución 3: el GEM, una red neuronal de grafos}
\label{sec:gem}
% 🔴 REESCRITO (sep-2026). El guion anterior defendia un Graph Transformer y «por que
%    atencion y no paso de mensajes». Es justo lo contrario de la decision vigente.
%    Fuente: doc/00_Justificaciones.md §5 — el cambio esta MEDIDO, no es de opinion.

\subsection{Por qué paso de mensajes y no atención}
% El argumento, que es una contribucion propia y no una cita:
%   - la alternativa de partida era el Graph Transformer de GTraQEM, que inyecta la
%     topologia mediante una MATRIZ DE ESTRUCTURA sumada al score de atencion;
%   - medido sobre nuestros grafos: la distancia mediana entre nodos conectados es de
%     30 a 70 pasos (p99 = 157), asi que con el recorte habitual esa matriz queda
%     CONSTANTE en mas del 99 % de sus entradas, y sin recorte no cabe en memoria;
%   - conclusion: a nuestra escala el mecanismo NO INFORMA;
%   - y GTraQEM ademas alimenta su modelo con el valor ruidoso medido, que este trabajo
%     excluye por diseño: sin esa entrada y sin matriz util, quedaba lo peor de ambos.

\subsection{El nodo virtual bidireccional}
% ⚠️ La BIDIRECCIONALIDAD es la diferencia clave con el diseño del paper, y hay que
%    explicarla: en un Transformer la atencion ya conecta todo con todo, asi que basta
%    un nodo virtual que LEA. En paso de mensajes no: los mensajes viajan un salto por
%    capa y solo por aristas, luego un nodo virtual sin salidas ACUMULA PERO NO DEVUELVE.
%    Con 4-6 capas frente a distancias de 30-70 pasos, la propagacion por aristas es
%    puramente local y el nodo virtual es la UNICA via al largo alcance del modelo.
% La idea del nodo virtual sí procede de GTraQEM; su implementacion, no.

\subsection{Los grafos completos, sin truncar}
% 🔴 ESTE APARTADO CAMBIA DE SIGNO. Antes planteaba una disyuntiva —truncar a 2 000
%    nodos, GPU grande o atencion eficiente— que YA NO EXISTE: con paso de mensajes la
%    memoria es LINEAL en nodos y aristas, no cuadratica.
% Lo que hay que contar es el efecto colateral, que es un argumento a favor:
%   desaparece el truncado y, con el, un sesgo metodologico que castigaba mas al
%   entrenamiento que al test.

% =====================================================================================
% EL PROTOCOLO IDENTICO
% =====================================================================================

\section{El objetivo común}
\label{sec:objetivo_comun}
Los tres modelos predicen $f$ y reportan sobre $\Delta$ reconstruido, según el
capítulo~\ref{cap:variable}. \textbf{Es lo que hace que la comparación sea justa.}

\section{La partición}
\label{sec:particion}
% Fuente: src/particion.py (modulo compartido por los tres modelos)
\subsection{Por qué la validación es temporal y no aleatoria}
El test es OOD en el tiempo: sus días de calibración son posteriores y disjuntos. Una
validación aleatoria compartiría días con el entrenamiento y daría una métrica
\textbf{optimista que no anticipa} lo que pasará en test.

Reservando los últimos días se obtiene una validación que \textbf{imita la relación real}
entre entrenamiento y test. Es más pequeña y más difícil, y esa es precisamente su utilidad.

\subsection{Una sola definición para los tres modelos}
La partición vive en un \textbf{único módulo} que importan tanto la vía tabular como el
cargador de grafos. \textbf{Es lo que convierte «protocolo idéntico» en algo demostrable} en
vez de en una afirmación.

Comprobación explícita: ningún día se comparte entre entrenamiento, validación y test.

\section{El preprocesado}
\label{sec:preprocesado}
% Fuente: las 17 directrices de notebooks/eda/04_directrices.ipynb seccion 4
Cada decisión sale de una medida del análisis exploratorio, no del criterio:
\begin{itemize}
    \item \textbf{normalización} obligatoria para Ridge (las variables abarcan diez órdenes de
          magnitud), innecesaria para Random Forest;
    \item \textbf{logaritmos}, porque la fidelidad cae de forma multiplicativa con el número
          de puertas;
    \item \textbf{poda} de las dimensiones constantes o duplicadas;
    \item \textbf{envoltura de ángulos}, porque una rotación y la misma más una vuelta
          completa son la misma operación física;
    \item \textbf{exclusión del día}, porque cae fuera del rango visto el $100$\,\% del tiempo.
\end{itemize}

\textbf{Regla que no se negocia:} todo se ajusta \textbf{solo con el conjunto de
entrenamiento}. Normalizar antes de particionar introduciría información del test.

\section{Métricas}
\label{sec:metricas}
MAE, RMSE, $R^2$ y mejora relativa frente a no mitigar.

\subsection{Cómo hay que reportarlas}
% Fuente: las 17 directrices. Varias salen de trampas detectadas.
\begin{itemize}
    \item \textbf{desglosadas por los tres ejes OOD} y por observable — un promedio global
          quedaría dominado por un solo tipo de circuito;
    \item \textbf{mediana y percentiles además de la media}, porque las distribuciones tienen
          colas muy pesadas;
    \item \textbf{también restringidas a las muestras con señal}, porque una fracción grande
          del conjunto no tiene nada que predecir e infla artificialmente cualquier promedio;
    \item en el eje de tamaño, \textbf{acompañadas de una métrica relativa}, porque el MAE baja
          solo porque el objetivo se encoge.
\end{itemize}

\subsection{El listón}
% ESTE es el criterio de exito del objetivo general del capitulo 5. Enlazarlos.
\textbf{No mitigar} es el baseline que da sentido al trabajo. Un modelo que no lo bata no
sirve, por bueno que sea su $R^2$.

\section{Reproducibilidad y trazabilidad}
\label{sec:mlops}
Semillas globales, versionado del conjunto de datos, registro de experimentos con MLflow, y el
manifiesto que acompaña al conjunto con versiones y configuración.

\section{Infraestructura de cómputo}
\label{sec:infra}
Generación en máquinas virtuales de Google Cloud (necesaria, no cómoda: las muestras de $15$
qubits piden $17{,}2$\,GB por proceso). Entrenamiento de la red de grafos en GPU.
